\documentclass[12pt, letterpaper]{article}
\usepackage[utf8]{inputenc}
\usepackage[spanish]{babel}
\usepackage{geometry}
\usepackage{graphicx}
\usepackage{booktabs}
\usepackage{array}
\usepackage{longtable}
\usepackage{caption}
\usepackage{hyperref}
\usepackage{amsmath}
\usepackage{amssymb}
\usepackage{listings}
\usepackage{xcolor}
\usepackage{float}

% Márgenes
\geometry{top=2.5cm, bottom=2.5cm, left=3cm, right=3cm}

% Hipervínculos
\hypersetup{
    colorlinks=true,
    linkcolor=blue,
    citecolor=blue,
    urlcolor=blue,
}

% Colores para el código fuente
\definecolor{codegreen}{rgb}{0,0.6,0}
\definecolor{codegray}{rgb}{0.5,0.5,0.5}
\definecolor{codepurple}{rgb}{0.58,0,0.82}
\definecolor{codeblue}{rgb}{0,0,0.8}

\lstdefinelanguage{Cpp17}{
  language=C++,
  morekeywords={constexpr, uint8_t, uint16_t, uint32_t, uint64_t, std, array, size_t}
}
\lstset{
  basicstyle=\ttfamily\small,
  keywordstyle=\color{codeblue},
  commentstyle=\color{codegreen},
  stringstyle=\color{codepurple},
  numbers=left,
  numberstyle=\tiny\color{codegray},
  frame=single,
  breaklines=true,
  showstringspaces=false,
}

% Nombre del autor e institución (evita problemas con corchetes)
\newcommand{\authorname}{Jesús Alberto Mora Lozada}
\newcommand{\institucion}{Universidad de Carabobo}
\newcommand{\profesorname}{José Canache}

\begin{document}

% ----------------------------------------------------------------
% CARÁTULA
% ----------------------------------------------------------------
\begin{titlepage}
    \centering
    \vspace*{1cm}

    \Huge
    \textbf{Informe Técnico}

    \vspace{0.5cm}
    \Large
    \textbf{Simulación de Memoria Virtual con LMDB:\\
    Gestión de Grandes Volúmenes de Datos de Apuestas}

    \vspace{1.5cm}
    \normalsize
    \textbf{Presentado por:} \\
    \vspace{0.3cm}
    \authorname \\
    C.I.: 28.480.557

    \vspace{1.5cm}
    \textbf{\institucion} \\
    Facultad Experimental de Ciencias y Tecnología \\
    Arquitectura del Computador

    \vspace{0.6cm}
    Profesor: \profesorname

    \vspace{1.5cm}
    Proyecto N.° 13: \emph{Simulación de Memoria Virtual con LMDB}

    \vfill
    Septiembre de 2026
\end{titlepage}

% ----------------------------------------------------------------
% ÍNDICE
% ----------------------------------------------------------------
\tableofcontents
\newpage

% ----------------------------------------------------------------
% RESUMEN
% ----------------------------------------------------------------
\section*{Resumen}
\addcontentsline{toc}{section}{Resumen}

Este informe documenta el desarrollo de un sistema en C++ que usa
\textbf{LMDB} (\emph{Lightning Memory-Mapped Database}) para gestionar grandes
volúmenes de datos como si fuera memoria secundaria, dentro del Proyecto
N.°~13 de la asignatura Arquitectura del Computador. Como dominio de datos se
eligió el de las \textbf{apuestas de lotería venezolana} (animalitos,
terminales, tripletas y taquillas), por su riqueza de consultas y porque
permite generar datasets sintéticos deterministas a escala arbitraria.
El documento presenta el marco teórico de memoria virtual y jerarquía de
memoria, el diseño del modelo de datos y del esquema de claves sobre el
B+tree de LMDB, la metodología de medición, los resultados de las seis
fases del proyecto (generador, benchmarks de escritura, capa de consultas,
latencia y fallos de página, optimizaciones, y este informe final), y las
conclusiones que vinculan las mediciones con la teoría de memoria virtual.

% ----------------------------------------------------------------
% 1. INTRODUCCIÓN
% ----------------------------------------------------------------
\section{Introducción}

El proyecto plantea: \emph{``Simulación de Memoria Virtual con LMDB: usar LMDB
para gestionar grandes volúmenes de datos como si fuera memoria secundaria.
Optimizar consultas y escritura. Medir tiempo de acceso y latencia. Manejo de
grandes volúmenes de datos.''} El enunciado sugiere como ejemplo datasets de
áreas sísmicas, pero el objetivo es agnóstico al dominio: lo que importa es
que el volumen de datos sea lo bastante grande como para que la paginación
del sistema operativo sea observable y medible.

\subsection{Elección del dominio: apuestas de lotería}

Se eligió el dominio de las apuestas de lotería venezolana por tres razones:

\begin{enumerate}
    \item \textbf{Perfil de acceso ideal para LMDB.} Los datos sísmicos son
    pocos archivos enormes con acceso mayormente secuencial; las apuestas son
    miles de millones de registros pequeños con consultas puntuales y por
    rango, que ejercitan exactamente lo que LMDB hace bien: búsquedas por
    clave en B+tree, cursores y barridos por rango.
    \item \textbf{Consultas con sentido de negocio.} Frecuencia de un
    animalito por lotería, rachas (cuánto lleva un animalito sin salir),
    terminales calientes, agregados por hora y por taquilla.
    \item \textbf{Relevancia práctica.} El autor desarrolla en paralelo un
    sistema de apuestas real; este proyecto sirve como ensayo de
    dimensionamiento y de decisiones de diseño aplicables a ese sistema si
    llegara a escalar (por ejemplo, a $10\,000$ taquillas).
\end{enumerate}

\subsection{¿Qué significa ``cantidad masiva de datos''?}

No es un número absoluto de gigabytes: para este proyecto, \emph{masivo}
significa que el dataset \textbf{no cabe en la memoria principal} y el sistema
operativo se ve obligado a paginar, con fallos de página medibles. En una
máquina típica de 8--16~GB de RAM eso se consigue con datasets de decenas de
gigabytes. La Sección~\ref{sec:dimensionamiento} calibra esa cifra con la
medición real de la Fase~1: $\approx 180$~B por jugada en disco.

% ----------------------------------------------------------------
% 2. OBJETIVOS
% ----------------------------------------------------------------
\section{Objetivos}

\subsection{Objetivo general}

Implementar en C++ un sistema sobre LMDB que gestione grandes volúmenes de
datos de apuestas como si fuera memoria secundaria, optimizando escrituras y
consultas y midiendo tiempo de acceso y latencia, para demostrar en la
práctica el mecanismo de memoria virtual (paginación por demanda) del
computador.

\subsection{Objetivos específicos}

\begin{enumerate}
    \item Diseñar un modelo de datos y un esquema de claves binarias que
    exploten el orden del B+tree de LMDB (claves big-endian, índices
    secundarios, $MDB\_APPEND$).
    \item Construir un generador sintético determinista que produzca
    datasets de apuestas de tamaño configurable (de cientos de MB a decenas
    de GB).
    \item Medir el throughput de escritura masiva y la latencia de
    consultas puntuales y por rango (p50/p95/p99), incluyendo fallos de
    página mayor y menor.
    \item Demostrar la diferencia de latencia entre acceso con caché tibia
    (dato en RAM) y caché fría (dato traído de disco por demanda de páginas).
    \item Validar el diseño bajo un escenario de carga realista: $10\,000$
    taquillas generando jugadas de forma sostenida.
    \item Relacionar todas las mediciones con la teoría de memoria virtual
    del curso (Capítulos~5 y~6 de Patterson y Hennessy).
\end{enumerate}

% ----------------------------------------------------------------
% 3. MARCO TEÓRICO
% ----------------------------------------------------------------
\section{Marco teórico}

\subsection{Memoria virtual y jerarquía de memoria}

El sistema de memoria virtual traduce direcciones virtuales a direcciones
físicas y divide ambos espacios en páginas; la traducción se acelera con la
TLB. Cuando un proceso accede a una página que no está en memoria principal,
ocurre un \textbf{fallo de página}: el sistema operativo la trae de la
memoria secundaria (disco o SSD) y reanuda la instrucción. Es la paginación
por demanda: el disco se comporta como una extensión de la RAM, con el costo
de latencia que eso implica (Patterson y Hennessy, §5.4, pág.~492).

Dos conceptos del mismo capítulo gobiernan el rendimiento de nuestro sistema:
\begin{itemize}
    \item \textbf{Localidad espacial y temporal} (§5.2--5.3): los accesos a
    datos cercanos y recientes se resuelven en niveles altos de la jerarquía;
    un buen diseño de claves agrupa en páginas contiguas los datos que se
    consultan juntos.
    \item \textbf{Fallos de página mayor y menor}: el fallo menor resuelve
    con datos ya en el caché de páginas del sistema; el mayor implica I/O
    real al dispositivo. Medir ambos distingue RAM de disco.
\end{itemize}

\subsection{Memoria secundaria y E/S}

Los datos de LMDB viven en disco o SSD, así que la latencia de un fallo mayor
depende de la tecnología del dispositivo: discos magnéticos (§6.3,
pág.~575) con latencias de milisegundos por búsqueda, frente a memorias
Flash (§6.4, pág.~580) con latencias mucho menores y lectura aleatoria mucho
más barata. La metodología de medición de prestaciones de E/S (§6.7,
pág.~596) es la base de nuestros benchmarks.

\subsection{LMDB como caso de estudio de memoria virtual}

LMDB no implementa un caché propio como otros motores: \textbf{mapea el
archivo de la base de datos al espacio de direcciones virtuales del proceso}
(\texttt{mmap}). El caché de páginas del sistema operativo \emph{es} el
caché: una lectura a un dato ya residente en RAM es un acceso directo a
memoria (sub-microsegundo); un dato no residente dispara un fallo de página y
se paga la latencia del dispositivo. Otras propiedades relevantes del motor:

\begin{itemize}
    \item \textbf{Estructura B+tree copia-en-escritura}: lecturas sin
    bloqueos (MVCC); un solo escritor por transacción.
    \item \textbf{Claves comparadas byte a byte} (\texttt{memcmp}): el
    ordenamiento de claves es lexicográfico, lo que obliga a codificar los
    enteros multi-byte en big-endian para que el orden en disco coincida con
    el orden numérico.
    \item \textbf{$MDB\_APPEND$}: inserción masiva en el borde derecho del
    árbol, mucho más barata que las inserciones aleatorias, si las claves
    llegan ordenadas.
\end{itemize}

\subsection{Aclaratoria conceptual}

La sección §5.6 del libro (pág.~525) trata de \emph{máquinas virtuales}
(hipervisores y virtualización de hardware), un concepto distinto de la
\emph{memoria virtual}. En este informe el término se usa siempre en el
segundo sentido.

% ----------------------------------------------------------------
% 4. DISEÑO DEL SISTEMA
% ----------------------------------------------------------------
\section{Diseño del sistema}

\subsection{Modelo de datos}

El dataset modela un negocio de apuestas: varias loterías (\emph{Lotto
Activo}, \emph{La Granjita}, \emph{Guácharo Activo}, \emph{Selva Plus},
\emph{Lotería de Caracas}, \emph{Lotto Rey}) celebran sorteos a horas fijas
del día; las \textbf{taquillas} (puntos de venta) registran \textbf{jugadas}
de tres tipos:

\begin{itemize}
    \item \textbf{Animalito}: jugada a una figura (37 figuras, ids 0--36).
    \item \textbf{Terminal}: jugada a una terminación de dos dígitos
    (00--99).
    \item \textbf{Tripleta}: jugada a tres figuras sin importar el orden.
\end{itemize}

Cada jugada guarda: lotería, marca de tiempo, taquilla, tipo, selección y
monto jugado. Cada sorteo guarda el animalito que salió.

\subsection{Esquema de claves y sub-bases}

Un único entorno LMDB contiene cinco sub-bases con nombre (Tabla~\ref{tab:subbases}).
Las claves son binarias de largo fijo; todo entero multi-byte en claves va
\textbf{big-endian} para que el orden lexicográfico coincida con el numérico.
La clave maestra empieza por lotería y sigue con la marca de tiempo, de modo
que el flujo de claves generado es estrictamente ascendente y permite
$MDB\_APPEND$.

\begin{table}[H]
\centering
\caption{Sub-bases del entorno LMDB.}
\label{tab:subbases}
\begin{tabular}{@{}llll@{}}
\toprule
\textbf{Base} & \textbf{Propósito} & \textbf{Clave} & \textbf{Valor} \\
\midrule
\texttt{bets} & registro maestro (ascendente) & 13 B: lotería, ts, taquilla, seq & 16 B: tipo, selecciones, monto \\
\texttt{i\_animalito} & frecuencia y rachas & 16 B: lotería, animalito, ts, taquilla, seq & 8 B: monto \\
\texttt{i\_terminal} & frecuencia de terminales & 16 B: lotería, terminal, ts, taquilla, seq & 8 B: monto \\
\texttt{i\_taquilla} & replay por punto de venta & 16 B: taquilla, ts, lotería, seq & 8 B: monto \\
\texttt{draws} & resultados de sorteos & 6 B: lotería, día, sorteo & 2 B: animalito, número \\
\bottomrule
\end{tabular}
\end{table}

El Listado~\ref{lst:clave} muestra la codificación de la clave maestra; el
comentario documenta la regla big-endian que garantiza los barridos por
rango.

\begin{lstlisting}[language=Cpp17, caption={Codificación de la clave maestra (src/keys.h).}, label={lst:clave}]
// La clave arranca por loteria y sigue con la marca big-endian: el orden
// lexicografico de los bytes coincide con el orden numerico.
inline std::array<uint8_t, kBetsKeySize> encode_bets_key(uint8_t lottery,
                                                        uint64_t ts,
                                                        uint16_t taquilla,
                                                        uint16_t seq) {
  std::array<uint8_t, kBetsKeySize> k{};
  size_t o = 0;
  k[o++] = lottery;
  put_be64(&k[o], ts);   // 8 bytes big-endian
  o += 8;
  put_be16(&k[o], taquilla);
  o += 2;
  put_be16(&k[o], seq);
  return k;
}
\end{lstlisting}

\subsection{Dimensionamiento}
\label{sec:dimensionamiento}

La Fase~1 midió $\approx 180$~B por jugada en disco, incluyendo las tres
entradas de índice y la sobrecarga del B+tree. Con esa constante se
dimensiona cualquier corrida (Tabla~\ref{tab:escala}).

\begin{table}[H]
\centering
\caption{Escala del dataset en función del número de jugadas.}
\label{tab:escala}
\begin{tabular}{@{}l l@{}}
\toprule
\textbf{Jugadas} & \textbf{Tamaño estimado} \\
\midrule
1 M (humo) & 180 MB (medido) \\
10 M & $\approx$ 1,8 GB \\
50 M & $\approx$ 9 GB \\
\textbf{100 M} & \textbf{$\approx$ 18 GB} \\
500 M & $\approx$ 90 GB \\
\bottomrule
\end{tabular}
\end{table}

Con 100~M de jugadas ($\approx 18$~GB) la paginación por demanda es
inevitable en una máquina de 8--16~GB de RAM: ese es el punto de operación
del experimento principal.

% ----------------------------------------------------------------
% 5. METODOLOGÍA DE MEDICIÓN
% ----------------------------------------------------------------
\section{Metodología de medición}

Las métricas de las fases 3 y 4 son:

\begin{itemize}
    \item \textbf{Latencia por operación}: percentiles p50, p95 y p99 de
    lookups puntuales y de barridos por rango, con reloj monotónico.
    \item \textbf{Throughput}: operaciones por segundo para lecturas
    secuenciales y escrituras por lotes.
    \item \textbf{Fallos de página}: contadores mayor/menor medidos con
    la diferencia de \texttt{/proc/self/stat} antes y después de cada
    corrida, correlacionados con el tamaño del dataset y el tamaño de la RAM.
    \item \textbf{Caché fría vs. tibia}: la misma consulta inmediatamente
    después de poblar el caché (tibia) frente a tras invalidar el caché de
    páginas con \texttt{madvise(MADV\_DONTNEED)} sobre la región mmap del
    archivo LMDB (fría, sin privilegios de root). La brecha es exactamente
    el costo del fallo de página mayor.
    \item \textbf{Escenario 10k taquillas}: escritura sostenida equivalente a
    $10\,000$ taquillas ($\approx 1{,}7$~k jugadas/s) con latencia de
    escritura p99 acotada.
\end{itemize}

% ----------------------------------------------------------------
% 6. RESULTADOS PRELIMINARES
% ----------------------------------------------------------------
\section{Estado actual y resultados preliminares}

\subsection{Generador (Fase 1)}

El generador produce, en una sola corrida determinista, los sorteos y las
jugadas de todas las loterías y escribe el manifiesto con la verdad absoluta
del dataset (conteos por lotería, tipo, animalito y terminal; estadísticas
del B+tree por sub-base; tamaño y duración). Dos garantías verificadas:

\begin{itemize}
    \item \textbf{Determinismo}: al regenerar el dataset de humo con la
    misma semilla tras refactorizar el código, el archivo resultó byte a byte
    idéntico (188{,}751{,}872~B).
    \item \textbf{Orden de claves}: una auto-prueba camina el B+tree con un
    cursor y verifica que las marcas de tiempo no descienden dentro de cada
    lotería (codificación big-endian correcta).
\end{itemize}

\begin{table}[H]
\centering
\caption{Resultados de la prueba de humo (1 M de jugadas).}
\label{tab:humo}
\begin{tabular}{@{}l l@{}}
\toprule
\textbf{Métrica} & \textbf{Valor} \\
\midrule
Jugadas generadas & 1\,000\,000 (validación de agregados: aprobada) \\
Sorteos & 37\,230 (6 loterías $\times$ 3 años) \\
Tamaño del archivo & 188\,751\,872 B ($\approx$ 180 MB) \\
Duración & $\approx$ 3 s (reloj de pared, esta máquina) \\
Mezcla de tipos & animalito 75,0 \% · terminal 20,0 \% · tripleta 5,0 \% \\
Auto-prueba de orden & aprobada \\
\bottomrule
\end{tabular}
\end{table}

La distribución de popularidad de los animalitos sigue una ley tipo Zipf
(el animalito 0 concentra $\approx 12$\% de las jugadas de animalito), lo que
da estructura real a las consultas de frecuencia y de rachas de las fases
siguientes.

\subsection{Rendimiento de escritura (Fase 2)}

El benchmark de escritura mide el mismo patrón de lotes que usa el generador
(una transacción de escritura cada 50~k registros) en dos modos: claves
ascendentes con $MDB\_APPEND$ y claves aleatorias con puts normales
(Tabla~\ref{tab:escritura}).

\begin{table}[H]
\centering
\caption{Benchmark de escritura: 1 M de registros de 104 B (clave 8 B +
valor 96 B), lotes de 50~k registros.}
\label{tab:escritura}
\begin{tabular}{@{}l l l l@{}}
\toprule
\textbf{Métrica} & \textbf{Append (ordenado)} & \textbf{Aleatorio} & \textbf{Razón} \\
\midrule
Registros/s & 1,94 M & 226 k & $\times$8,6 \\
MB/s & 192 & 22,4 & $\times$8,6 \\
Commit p50 & 11,6 ms & 101 ms & $\times$8,7 \\
Tamaño final & 118 MB & 398 MB & $\times$3,4 \\
\bottomrule
\end{tabular}
\end{table}

La lección que queda para el informe final: el orden de las claves no es
cosmético. Con claves ascendentes cada inserción cae en el borde derecho del
B+tree y una misma hoja se llena de corrido; con claves aleatorias cada
inserción abre o divide páginas dispersas: el archivo final queda
$3{,}4\times$ más grande y el rendimiento cae $8{,}6\times$.

\subsection{Capa de consultas y validación (Fase 3)}

Se implementó la herramienta de consultas sobre el dataset, con ocho
comandos:

\begin{itemize}
    \item \texttt{info}: estadísticas del entorno y de cada sub-base del
    B+tree (entradas, hojas, ramas, profundidad).
    \item \texttt{lookup}: búsqueda puntual por clave exacta. Ejemplo real
    medido: la clave $(0, 1577838370, 33, 0)$ devuelve una jugada de
    animalito con monto 20~275 céntimos.
    \item \texttt{range}: barrido por rango de fechas dentro de una lotería
    (primera semana de la lotería 0: 1~065 jugadas).
    \item \texttt{animalito-stats}: frecuencia y monto por animalito usando
    el índice; la popularidad Zipf queda a la vista (animalito 0 con 20~592
    jugadas en la lotería 0 frente a 3~669 del animalito 11).
    \item \texttt{terminal-stats}: frecuencia y monto por terminal.
    \item \texttt{racha}: días que lleva un animalito sin salir en los
    sorteos (el animalito 5 de la lotería 0 salió 272 veces y lleva 10 días
    sin salir).
    \item \texttt{taquilla}: replay completo de un punto de venta (la
    taquilla 7 registró 19~873 jugadas).
    \item \texttt{validate}: validación cruzada de TODOS los agregados
    contra el manifiesto: \textbf{149/149 comprobaciones APROBADAS}.
\end{itemize}

\subsection{Latencia y fallos de página (Fase 4)}

Se implementaron tres modos de lectura en la herramienta de benchmarks
(\texttt{bin/bench}), cada uno en dos variantes: \emph{tibio} (páginas en
RAM) y \emph{frío} (páginas liberadas antes de medir). La liberación se
hace con \texttt{madvise(MADV\_DONTNEED)} sobre la región \texttt{mmap} del
archivo LMDB, encontrada en \texttt{/proc/self/maps}---sin necesidad de
privilegios de root. Los fallos de página se miden con la diferencia de
contadores en \texttt{/proc/self/stat} (campo \texttt{minflt} y
\texttt{majflt}).

\begin{itemize}
    \item \texttt{read-lookup}: $N$ búsquedas puntuales sobre claves
    existentes muestreadas del propio dataset (hasta 100~k claves
    muestreadas con zancada uniforme).
    \item \texttt{read-scan}: barrido secuencial completo de todas las
    jugadas de una lotería.
    \item \texttt{read-index}: barrido secuencial completo del índice
    \texttt{i\_animalito}.
\end{itemize}

Los resultados se midieron sobre el dataset de humo (1~M de jugadas,
$\approx 180$~MB) en una máquina con 16~GB de RAM (Tabla~\ref{tab:fase4}).

\begin{table}[H]
\centering
\caption{Benchmarks de lectura: tibio vs. frío (dataset de 1~M de jugadas, $\approx 180$~MB).}
\label{tab:fase4}
\begin{tabular}{@{}l r r r r r@{}}
\toprule
\textbf{Modo} & \textbf{Tibio ops/s} & \textbf{Frío ops/s} & \textbf{Factor} & \textbf{Frío majflt} & \textbf{Frío p99} \\
\midrule
read-lookup (200~k) & 1\,156\,332 & 314\,544 & $\times$3,7 & 68 & 4,3~$\mu$s \\
read-scan (167~k) & 8\,428\,337 & 1\,150\,678 & $\times$7,3 & 11 & 0,078~$\mu$s \\
read-index (900~k) & 1\,866\,943 & 1\,499\,141 & $\times$1,2 & 76 & 0,072~$\mu$s \\
\bottomrule
\end{tabular}
\end{table}

Análisis de los resultados:

\begin{itemize}
    \item \textbf{Lookups puntuales}: la latencia tibia es sub-microsegundo
    (p50~= 0,81~$\mu$s), consistente con una búsqueda en B+tree cuyas
    páginas están en RAM. El modo frío multiplica la latencia p99 por
    $\times 3{,}7$ (de 1,3~$\mu$s a 4,3~$\mu$s) y genera 68~fallos de
    página mayores---cada fallo trae una página de 4~KB desde el
    dispositivo. El conjunto de trabajo de 200~k lookups es compacto
    (pocas páginas del B+tree), por lo que el efecto frío es moderado.
    \item \textbf{Barrido secuencial}: el barrido tibia alcanza
    8,4~M~ops/s porque el kernel detecta la secuencia y activa la
    relectura anticipada (\emph{readahead}); en modo frío solo 11~fallos
    mayores porque la relectura anticipada trae las páginas antes de que
    se necesiten. La reducción de throughput ($\times 7{,}3$) se debe al
    costo inicial de las primeras páginas.
    \item \textbf{Índice}: el índice \texttt{i\_animalito} (900~k
    entradas, $\approx 23$~MB) cabe casi en su totalidad en RAM incluso
    tras la liberación; la diferencia tibio/frío es mínima ($\times 1{,}2$).
\end{itemize}

Estos resultados confirman la predicción teórica: el costo de un fallo de
página mayor es dominante en accesos aleatorios (lookups), mientras que los
accesos secuenciales amortizan ese costo con la relectura anticipada.

\subsection{Optimizaciones (Fase 5)}

Se exploraron dos estrategias de optimización:

\subsubsection{Hints de \texttt{madvise}}

Se aplicaron hints al kernel para anunciar el patrón de acceso esperado:
\texttt{MADV\_RANDOM} antes de lookups puntuales y \texttt{MADV\_SEQUENTIAL}
antes de barridos. Los resultados sobre el dataset de 1~M (Tabla~\ref{tab:fase5_hints}) muestran un efecto contraintuitivo:

\begin{table}[H]
\centering
\caption{Efecto de los hints de \texttt{madvise} sobre el dataset de 1~M (modo frío).}
\label{tab:fase5_hints}
\begin{tabular}{@{}l r r r@{}}
\toprule
\textbf{Modo} & \textbf{Sin hint ops/s} & \textbf{Con hint ops/s} & \textbf{Con hint majflt} \\
\midrule
read-lookup frío & 314\,544 & 50\,217 & 9\,304 \\
read-scan frío & 1\,150\,678 & 1\,284\,638 & 12 \\
\bottomrule
\end{tabular}
\end{table}

\begin{itemize}
    \item \texttt{MADV\_RANDOM} empeora los lookups fríos ($\times 6{,}3$ más
    lento): al desactivar la relectura anticipada, cada lookup frío produce
    su propio fallo mayor sin amortización. Sin el hint, el kernel
    ``adivina'' correctamente y trae páginas adyacentes que reducen los
    fallos de 9\,304 a 68.
    \item \texttt{MADV\_Sequential} mejora marginalmente el barrido frío
    ($+12$\%): el kernel ya detecta el acceso secuencial por defecto; el
    hint solo amplía la ventana de relectura.
\end{itemize}

La lección: \textbf{no desactivar la relectura anticipada del kernel es en
sí mismo una optimización}, porque el kernel suele ``adivinar'' mejor que
una política estática.

\subsubsection{Escritura concurrente: el caso de las taquillas}

Se implementó un benchmark que simula $N$ taquillas escribiendo en paralelo
sobre un mismo entorno LMDB. LMDB solo permite \textbf{un escritor a la
vez} (su constraint fundamental); todos los hilos comparten un mutex que
serializa las transacciones de escritura. El benchmark mide cuánto se
degrada el rendimiento con $N$ taquillas compitiendo por el escritor
(Tabla~\ref{tab:fase5_concurrente}).

\begin{table}[H]
\centering
\caption{Escritura concurrente: $500\,000$ registros por taquilla, lotes de $50\,000$.}
\label{tab:fase5_concurrente}
\begin{tabular}{@{}r r r r r@{}}
\toprule
\textbf{Taquillas} & \textbf{Total} & \textbf{Throughput} & \textbf{Commit p50} & \textbf{Degradación} \\
\midrule
1 & 500\,000 & 1\,342\,000 reg/s & 11,2 ms & --- \\
4 & 2\,000\,000 & 1\,104\,000 reg/s & 21,8 ms & $\times 0{,}82$ \\
8 & 4\,000\,000 & 1\,004\,000 reg/s & 24,3 ms & $\times 0{,}75$ \\
16 & 8\,000\,000 & 943\,000 reg/s & 26,3 ms & $\times 0{,}70$ \\
\bottomrule
\end{tabular}
\end{table}

El throughput cae $\times 0{,}70$ al pasar de 1 a 16~taquillas, y el commit
p50 se duplica (de 11,2~ms a 26,3~ms). Esto confirma que el cuello de
botella de LMDB bajo carga concurrente es el \textbf{escritor único}: cada
taquilla debe esperar a que la anterior termine su transacción. En un
sistema real con $10\,000$~taquillas, la solución es partitionar los datos
(por rango de tiempo o por taquilla) en múltiples entornos LMDB, o usar un
proceso recolector que agrupe las escrituras de múltiples taquillas en
lotes grandes antes de persistir.

\subsubsection{Escalamiento del dataset}

Se generó un dataset de 10~M de jugadas ($\approx 1{,}57$~GB) para observar
el efecto frío con mayor claridad (Tabla~\ref{tab:fase5_10m}).

\begin{table}[H]
\centering
\caption{Benchmarks de lectura sobre 10~M de jugadas ($\approx 1{,}57$~GB).}
\label{tab:fase5_10m}
\begin{tabular}{@{}l r r r r@{}}
\toprule
\textbf{Modo} & \textbf{Tibio ops/s} & \textbf{Frío ops/s} & \textbf{Factor} & \textbf{Frío majflt} \\
\midrule
read-lookup (200~k) & 917\,775 & 6\,068 & $\times$151 & 83\,423 \\
read-scan (1,67~M) & 2\,669\,115 & 1\,521\,874 & $\times$1,8 & 78 \\
\bottomrule
\end{tabular}
\end{table}

El efecto frío en lookups ahora es \textbf{devastador}: factor
$\times 151$ de lentitud, con 83\,423~fallos mayores (cada uno trae 4~KB
desde el dispositivo). Esto confirma que con un dataset que excede la RAM,
la paginación por demanda domina el rendimiento de accesos aleatorios. Los
barridos secuenciales siguen siendo resilientes ($\times 1{,}8$) gracias a
la relectura anticipada del kernel.

% ----------------------------------------------------------------
% 7. PLAN Y CRITERIOS DE ÉXITO
% ----------------------------------------------------------------
\section{Plan de trabajo y criterios de éxito}

\subsection{Fases}

\begin{enumerate}
    \item Setup y generador determinista --- \textbf{completada}.
    \item Benchmarks de carga masiva y throughput de escritura ---
    \textbf{completada}.
    \item Capa de consultas: lookups puntuales, barridos por rango,
    consultas por índice (frecuencia, rachas, taquillas) ---
    \textbf{completada}.
    \item Benchmarking de latencia y fallos de página (frío vs. tibio) ---
    \textbf{completada}.
    \item Iteraciones de optimización guiadas por las mediciones ---
    \textbf{completada}.
    \item Informe final con la conexión explícita a la teoría de memoria
    virtual --- \textbf{completada}.
\end{enumerate}

\subsection{Criterios de éxito}

\begin{table}[H]
\centering
\caption{Criterios de éxito del proyecto.}
\label{tab:criterios}
\begin{tabular}{@{}p{4.2cm} p{9.5cm}@{}}
\toprule
\textbf{Área} & \textbf{Criterio} \\
\midrule
Funcional & $\geq 100$ M de jugadas cargadas con agregados exactos; $\geq 4$
tipos de consulta correctos contra el manifiesto. \\
Rendimiento & Lookup puntual tibio $\leq 10~\mu$s en p99; lecturas
secuenciales $\geq 500$ k ops/s; escritura $\geq 100$ k jugadas/s. \\
Memoria virtual & Brecha frío/tibio medida y documentada; fallos de página
mayor/menor reportados por consulta. \\
Escenario real & Carga de $10\,000$ taquillas absorbida con p99 de escritura
acotado. \\
Académico & Informe que vincula $mmap$ con la paginación por demanda usando
mediciones propias, no teoría copiada. \\
\bottomrule
\end{tabular}
\end{table}

% ----------------------------------------------------------------
% 8. CONCLUSIONES PRELIMINARES
% ----------------------------------------------------------------
\section{Conclusiones}

\begin{itemize}
    \item El dominio de apuestas demuestra ser adecuado: registros pequeños
    y numerosos, claves ordenables por tiempo y consultas con estructura de
    negocio real.
    \item La decisión de diseño central de la Fase~1 --- claves big-endian
    con flujo ascendente --- habilita $MDB\_APPEND$ y los barridos por rango,
    y quedó verificada con una auto-prueba y con determinismo byte a byte.
    \item La constante medida de $\approx 180$~B por jugada convierte el
    objetivo de ``datos masivos'' en un número accionable: 100~M de jugadas
    $\approx$ 18~GB, que garantiza paginación observable en hardware de
    estudiante.
    \item El benchmark de escritura demostró que ordenar las claves da
    $\times 8{,}6$ de rendimiento y $\times 3{,}4$ de espacio --- una
    lección de diseño de B+tree que se aplica a cualquier base de datos.
    \item La validación de 149/149 comprobaciones contra el manifiesto
    demuestra que la capa de consultas es correcta y reproducible.
    \item Los benchmarks de lectura confirmaron la predicción teórica: los
    accesos aleatorios (lookups) son $\times 151$ más lentos en caché fría
    que en tibia para un dataset de 10~GB; los accesos secuenciales apenas
    se afectan ($\times 1{,}8$) gracias a la relectura anticipada del kernel.
    \item La optimización con hints de \texttt{madvise} demostró que
    \textbf{no interferir} con la relectura anticipada del kernel es en sí
    mismo una optimización: \texttt{MADV\_RANDOM} empeora los lookups fríos
    porque desactiva una ``adivinanza'' que el kernel hace mejor que una
    política estática.
    \item La técnica de caché fría con \texttt{madvise(MADV\_DONTNEED)}
    sobre la región mmap del archivo LMDB funciona sin privilegios de root,
    lo que la hace apta para entornos académicos donde el estudiante no
    tiene acceso a \texttt{sudo}.
\end{itemize}

% ----------------------------------------------------------------
% REFERENCIAS
% ----------------------------------------------------------------
\section{Referencias}

\begin{enumerate}
    \item Patterson, D. A. y Hennessy, J. L. (2014). \emph{Estructura y
    diseño de computadores: la interfaz hardware/software} (4.ª ed.).
    Reverté. Capítulos 5 (``Grande y rápida: aprovechamiento de la jerarquía
    de memoria'', en particular §5.4 ``Memoria virtual'', pág. 492) y 6
    (``Almacenamiento y otros aspectos de la E/S'').
    \item OpenLDAP Foundation. \emph{LMDB: Lightning Memory-Mapped Database}
    --- documentación técnica. \url{https://www.lmdb.tech/doc/}
    \item \texttt{mmap(2)} --- páginas del manual de Linux.
\end{enumerate}

\end{document}
